\documentclass[tikz,border=8pt]{standalone}
\usepackage{xcolor}
\definecolor{bg}{HTML}{171D2B}
\definecolor{aksen}{HTML}{F7B538}
\begin{document}
\pagecolor{bg}
\begin{tikzpicture}[
  x=1.15cm, y=1.15cm,
  line cap=round, line join=round,
  titik/.style={circle, fill=aksen, inner sep=2.4pt},
  label putih/.style={text=white, font=\large},
]
  % Data: frekuensi harian menabung (Rupiah): 1000,2000,2000,3000,5000,5000,5000,10000
  % Nilai kumulatif: 1000:1, 2000:2, 3000:1, 5000:3, 10000:1
  % Sumbu x: nominal (1000..10000), sumbu y: frekuensi (0..3)
  \draw[white!70, thick, ->] (0,0) -- (10.8,0);
  \draw[white!70, thick, ->] (0,0) -- (0,4.2);

  % garis bantu horizontal
  \foreach \y in {1,2,3} {
    \draw[white!25, dashed, thin] (0,\y) -- (10.2,\y);
  }

  % label sumbu y
  \foreach \y in {1,2,3} {
    \node[label putih, left=4pt] at (0,\y) {\y};
  }
  \node[label putih, rotate=90, above=22pt] at (0,2.2) {Frekuensi};

  % titik-titik data: (1000,1) (2000,2) (3000,1) (5000,3) (10000,1)
  % skala x: 1 unit = 1000 rupiah
  \draw[aksen, very thick] (1,1) -- (2,2) -- (3,1) -- (5,3) -- (10,1);
  \node[titik] at (1,1) {};
  \node[titik] at (2,2) {};
  \node[titik] at (3,1) {};
  \node[titik] at (5,3) {};
  \node[titik] at (10,1) {};

  % label sumbu x (nominal)
  \node[label putih, below=6pt] at (1,0) {1000};
  \node[label putih, below=6pt] at (2,0) {2000};
  \node[label putih, below=6pt] at (3,0) {3000};
  \node[label putih, below=6pt] at (5,0) {5000};
  \node[label putih, below=6pt] at (10,0) {10000};
  \node[text=aksen, font=\large\bfseries, below left=2pt] at (0,0) {0};
  \node[label putih, below=24pt] at (5.5,0) {Nominal uang jajan (Rupiah)};
\end{tikzpicture}
\end{document}
